\section{A numerical triple-genus bound for every alternating or symmetric cover}
\label{sec:effective-triple-genus}

\paragraph{Status.}
HD-GC03--07 give a new proof of an explicit numerical form of the
known qualitative triple-genus growth. The proof was written by root
on 22 November 2026. Complete author and separate independent whole-proof readings pass. Root has also read the full independent audit, its entire exact checker and first successful output, and both targeted derivation reports. The two recorded wording repairs change no mathematical display.
The cycle-mode argument is a numerical specialization of the method
in Neftin--Zieve, Proposition~7.1 and Corollary~7.3; all constants
and the complete argument needed here are proved below. No statement
from that preprint with unspecified constants is used as a premise.
HD-GC01--02 and the standard finite automorphism description of a
complex elliptic curve are the previously checked inputs.

Use the actual covers and notation of HD-GC01--02:
\[
 g=g(Y_1),\qquad h=g(Y_2),\qquad b=g(X_2),\qquad c=g(X_3),
 \qquad A=R_{Y_2/Y_1}.
\]
The base is $\mathbf P^1$, the Galois group is $A_n$ or $S_n$, and
all quotient curves and maps are the actual connected ones. In this
section $n\geq10^6$. Identity inertia may always be omitted.

\subsection*{HD-GC03: an explicit cycle-mode lemma}

Let a permutation on $n$ letters have individual cycle lengths $r_i$,
and write $A(x)=\sum_{i,j}(r_i-\gcd(r_i,r_j))$ as in HD-GC01.
If $A(x)<4n$, either there is a $k\in\{1,\ldots,6\}$ such that
all cycles other than the $k$-cycles have total length at most
$4k(k+1)$, or each cycle length $k\leq6$ occurs at most eight times.
More generally, if $A(x)<6n$, the corresponding alternatives hold
with $6k(k+1)$ and twelve times, respectively.

\begin{proof}
We prove both statements simultaneously. Let $q=8$ in the first
case and $q=12$ in the second, so $A(x)<qn/2$. If some length at
most six has more than $q$ cycles, let $k$ be the least such length
and let $m>q$ be its multiplicity. The sum $R$ of lengths smaller
than $k$ is at most $qk(k-1)/2$. Comparing the $m$ cycles of length
$k$ with cycles of length greater than $k$ gives
\[
       \frac m2(n-km-R)\leq A(x)<\frac q2 n,
       \qquad m(n-km-R)<qn.                         \tag{EG1}
\]
Here $r-\gcd(r,k)\geq r/2$ for every $r>k$.

First suppose $m\leq n/(2k)$. Since $R\leq180$, (EG1) gives
$m<qn/(n/2-180)<2q+1$, hence $m\leq2q$. Then
$km+R\leq144+180=324$ gives
$m<qn/(n-324)<q+1$, a contradiction. All the strict numerical
comparisons hold for both $q$ at $n\geq10^6$.
Thus $m>n/(2k)$. Put $E=n-km$. Then
\[
 E-R<2qk,\qquad
 E<R+2qk\leq\frac q2 k(k+3)\leq324.
\]
Using $m=(n-E)/k$ once more in (EG1) gives
$E-R<qkn/(n-324)<qk+1$. Since $E-R$ is an integer,
\[
             E\leq R+qk\leq\frac q2 k(k+1).
\]
If no such least length exists, the second alternative holds by
definition. This proves both versions.
\end{proof}

Assign a permutation its finite mode $k$ by the construction just
given, and assign mode $\infty$ in the second alternative. For the
$4n$ version, a finite mode has exceptional length $E\leq168$;
mode one moves at most eight letters. A mode-$\infty$ permutation
has at most $48$ cycles of length at most six, and therefore at most
$n/7+48$ cycles in total. Consequently its natural index is at least
$6n/7-48$. For finite $k>1$, its index differs from $n(1-1/k)$ by
at most $168$, and is at least $n/2-168$.

\subsection*{HD-GC04: the low ordered-pair-genus case}

Assume $g\leq1$ and $h<n$. Then
\[
        A=2(h-1)-2(n-1)(g-1)<4n.                    \tag{EG2}
\]
Every local $A(x)$ is nonnegative, so HD-GC03 applies to every
branch point. Call modes greater than one large, including mode
$\infty$, and call mode one small. Let $J$ be the sum of the natural
indices of the small branch permutations.

For a small permutation with moved support $E\leq8$ and natural
index $j$, its fixed/nontrivial cycle pairs contribute exactly
$(n-E)j$ to $A(x)$. The remaining contribution is nonnegative.
Thus
\[
       A\geq(n-8)J,\qquad J<\frac{4n}{n-8}<5,
       \qquad J\leq4.                              \tag{EG3}
\]
There are at most four large branches: five would have total natural
index at least $5(n/2-168)>2n$, whereas natural Riemann--Hurwitz gives
total index $2n+2g-2\leq2n$.

The multiset of large modes is one of
\[
 (3,3,3),\ (2,4,4),\ (2,3,6),\ (2,2,2,2),\
 (\infty,2,2),\ (\infty,\infty).                    \tag{EG4}
\]
Here and below all constants are numerical. To verify the list, first
suppose every large mode is finite, and let
$S=\sum(1-1/k)$. Natural Riemann--Hurwitz and (EG3) give
$|S-2|\leq678/n<1/60$. As $S$ has denominator dividing $60$, $S=2$.
For three terms the reciprocal sum is one, giving exactly
$(3,3,3),(2,4,4),(2,3,6)$; four terms must all be two. At most two
finite terms cannot sum to two.

If there is one infinite mode and the finite modes have weight sum
$S$, its index is $(2-S)n+\epsilon$ with $|\epsilon|\leq510$.
The index bounds $6n/7-48\leq I\leq n$ give
\[
           1-510/n\leq S\leq8/7+558/n.
\]
One finite term is at most $5/6$, three are at least $3/2$,
and two terms are either $1$ or at least $7/6$.
Since $7/6-8/7=1/42>558/n$, the only possibility is two modes two.
Two infinite modes cannot be accompanied by a finite mode, since
$2(6n/7-48)+(n/2-168)>2n$; three infinite modes are also impossible.
This proves (EG4).

We claim that, apart from the branch pattern in HD-GC05 below,
\[
                              h\geq n/24.           \tag{EG5}
\]
There are three cases.

\paragraph{A finite even mode.}
Such a mode has $k\in\{2,4,6\}$ and contains at least
$n/k-4(k+1)\geq n/6-28$ even cycles. Each even cycle contributes
one to the ramification of $Y_2\to X_2$: its antipodal pairs form
one unordered-pair orbit whose ordered double cover is ramified.
All other pair orbits are unramified in that double cover. Hence
\[
       h=2b-1+\tfrac12R_{\pi_2}
          \geq n/12-15\geq n/24.                   \tag{EG6}
\]
This treats every list in (EG4) except $(3,3,3)$ and
$(\infty,\infty)$.

\paragraph{Three modes three.}
For each of the three permutations let its exceptional cycles have
total length $E_i\leq48$, number $k_i$, and number $N_i$ whose
length is not divisible by three. Put
\[
 S_i=\sum_{\rm exceptional}(r-3),\qquad
 Q_i=S_i+4N_i,
 \qquad N=N_1+N_2+N_3.
\]
Pairing an exceptional cycle of length $r$ with a main three-cycle
in both orders contributes $r+3-2\gcd(r,3)$, which is $r-3$ if
$3\mid r$ and $r+1$ otherwise. Thus its exact pair-ramification
contribution is
\[
 A_i=\frac{n-E_i}{3}Q_i+A_{{\rm exc},i},
 \qquad 0\leq Q_i\leq2E_i,\quad
 0\leq A_{{\rm exc},i}\leq E_i^2.
\]
It follows that $A_i-nQ_i/3$ lies between $-1536$ and $2304$.
For a small branch with moved support $E\leq8$, number $k$ of
nontrivial cycles and index $j$, the remaining internal contribution
is at most $(k-1)j\leq Ej$. Thus its total contribution differs
from $nj$ by a number in $[-8j,0]$. Summing and using (EG3),
the total error lies in $[-4640,6912]$, and in particular in the
slightly wider interval $[-9248,7168]$ used below.

The natural indices of the three large permutations are
$2n/3+S_i/3$. Therefore natural Riemann--Hurwitz gives
$\sum_i S_i+3J=6g-6$. The $nJ$ term from the small branches
cancels exactly, leaving
\[
 \begin{split}
 A&=n(2g-2+4N/3)+\epsilon,
           \qquad -9248\leq\epsilon\leq7168,\\
 h&=n(2g-2+2N/3)+(2-g)+\epsilon/2.                 \tag{EG7}
 \end{split}
\]
If $g=0$ and $N\leq2$, its upper bound makes $h<0$, impossible.
If $g=0,N=3$, the equality $\sum_iS_i+3J=-6$ forces exactly three
exceptional fixed-point cycles, no other exceptional cycles, and
$J=0$: each of these three non-three-divisible cycles contributes
at least $-2$, and all remaining summands are nonnegative.
If $g=1,N=0$, the same equality forces no exceptional cycles and
$J=0$. In either case all three genuine branch permutations have
order three. The actual Galois-cover Riemann--Hurwitz formula then
gives genus one for the Galois cover. A finite group of automorphisms
of a complex genus-one curve is an extension of an abelian
translation group by a cyclic group, hence solvable, contradicting
$A_n$ or $S_n$ for $n\geq5$.
\cite[Chapter IV, Proposition 7.13 and Remark 7.18, pp.~140--141]{TD:MilneEC2021}
Thus $N\geq4$ if $g=0$, and $N\geq1$ if $g=1$. Equation (EG7)
gives $h\geq(2/3)n-4623\geq n/2$, proving (EG5) here.

\paragraph{Two infinite modes.}
Let $u,v$ be the numbers of cycles in the two large permutations.
The exact natural Riemann--Hurwitz identity is
\[
                         J=u+v+2g-2.               \tag{EG8}
\]
If $g=1$, then $J\geq2$, so
$h=1+A/2\geq1+(n-8)J/2\geq n-7$.
If $g=0$ and $J\geq3$, then
$h=2-n+A/2\geq2-n+(n-8)J/2\geq n/2-10$.
Both imply (EG5).

It remains that $g=0$ and $J\leq2$. Since $u,v\geq1$, if $J=0$
then $u=v=1$ and $A=0$, giving $h=2-n<0$; this is impossible.
If $J=1$, the two cycle numbers are one and two, and the small
branch data consist of one transposition. This is exactly the
special pattern treated in HD-GC05.

If $J=2$, then $u+v=4$. An unequal partition into $u$ lengths
has $A(x)\geq n/(2u)$: compare a largest and a smallest length,
whose gcd is at most half the largest. Thus if either of the two
large partitions is unequal, $u,v\leq3$ and
\[
       h=2-n+A/2\geq2-n+(2(n-8)+n/6)/2
                         =n/12-6\geq n/24.
\]
If both partitions are equal-length partitions, their lengths are
$n/u$ and $n/v$. Here $(u,v)$ is $(1,3),(2,2)$ or $(3,1)$.
Choose a coordinate on the genus-zero natural curve and send the
two corresponding base values to zero and infinity. The defining
rational function has every zero and pole multiplicity divisible
by $d=n/\operatorname{lcm}(u,v)\geq n/3$. Factoring its zeros and
poles over $\mathbf C$ writes it as a $d$-th power of a rational
function of degree $\operatorname{lcm}(u,v)\in\{2,3\}$.
Both factors have degree greater than one. Their intermediate
fibres give a nontrivial block system on the $n$ letters,
contradicting the primitive natural action of $A_n$ or $S_n$.
This completes the proof of (EG5), with only the stated special
pattern excluded from that intermediate conclusion.

\subsection*{HD-GC05: the exceptional three-branch pattern}

Suppose the natural genus is zero and the branch cycles have types
\[
                   [n],\quad[a,n-a],\quad[2,1^{n-2}].
\]
Put $a'=n-a$ and $d=\gcd(a,a')$. The pair ramification degrees
of these three permutations are $0,n-2d,n-2$, respectively.
Therefore $h=1-d$. Connectedness of the actual ordered-pair
quotient forces $d=1$ and $h=0$.

Put $D=n(n-1)(n-2)$. The $n$-cycle has $(n-1)(n-2)$ orbits on
ordered distinct triples. The split permutation has
\[
 (a-1)(a-2)+(a'-1)(a'-2)+3(a-1)+3(a'-1)
                         =a^2+(a')^2-2
\]
such orbits: mixed triples have orbit length $aa'$ because
$\gcd(a,a')=1$. The transposition has index $3(n-2)^2$ on
ordered triples, obtained by subtracting its
$(n-2)(n-3)(n-4)$ fixed triples from $D$ and dividing by two.
Riemann--Hurwitz consequently gives the exact identity
\[
 2(g(Y_3)-1)=(n-5)(n-2)+2(a-1)(a'-1).              \tag{EG9}
\]
The local formula HD-GC01 bounds the three contributions to
$R_{\pi_3}$ by $4,11,3(n-2)$, respectively. In the middle term
there is at most one ordered pair with unequal two-adic valuations,
whose gcd is one, and at most two cycles divisible by three.
Thus $R_{\pi_3}\leq3n+9$. From (EG9) and the degree-six map
$Y_3\to X_3$ we obtain
\[
                  12c\geq n^2-10n+13>n^2/2
                         \qquad(n\geq20).          \tag{EG10}
\]
No classification theorem or realization of arbitrary branch lists
is inferred: this is a calculation for an actual cover with this list.

\subsection*{HD-GC06: the uniform numerical result}

For every actual connected $A_n$- or $S_n$-Galois cover of the
sphere and $n\geq10^6$, its unordered-triple quotient satisfies
\[
                         \boxed{\ g(X_3)>n^2/288\ }. \tag{EG11}
\]
Indeed HD-GC02 proves the stronger $c>n^2/12$ if $g\geq2$ or
$h\geq n$. Otherwise HD-GC04 gives $h\geq n/24$, apart from
the special case settled more strongly by (EG10). Substituting
$h\geq n/24$ and $g\geq0$ in (GC1) yields
\[
 12(c-b)\geq n^2/12-(38/3)n+26>n^2/24,
\]
where the last inequality holds already for $n\geq304$.
Since $b\geq0$, (EG11) follows.

In particular every fixed genus bound $G\geq0$ is exceeded by
the triple quotient once
\[
                 n\geq\max\{10^6,\lceil\sqrt{288G}\rceil\}.
\]
This supplies the numerical triple-genus input in HD-FT07--12:
one can replace their symbolic bound $g(X_3)>c_0n^2/12$ by
$c_0=1/24$ with threshold $10^6$. It does not supply a numerical
version of every subgroup or pair-genus classification theorem.

\subsection*{HD-GC07: an effective branch-count alternative for tail two}

Suppose the unordered-pair quotient has $b\leq14$. Then, for
$n\geq10^6$, the natural quotient has at most eight genuine
branch points. This consequence does not use $N_{14}$.

First $g\leq b$. To see this directly, the permutation module on
letters embeds in the module on unordered pairs by sending a
letter $i$ to the sum of pairs containing it. The map is injective:
vanishing of all pair coefficients means $a_i+a_j=0$ for all
$i\ne j$, whence every $a_i=0$ for $n\geq3$.
Taking multiplicities in the actual Galois cover's holomorphic
differentials identifies the two invariant dimensions with $g$
and $b$, respectively. Semisimplicity over $\mathbf C$ therefore
gives $g\leq b\leq14$.

Every even cycle contributes one to $R_{\pi_2}$ and at least one
to the natural index, so $R_{\pi_2}\leq2n+2g-2$. Consequently
\[
 h=2b-1+R_{\pi_2}/2\leq n+40,\qquad
 A=2(h-1)-2(n-1)(g-1)\leq4n+76<6n.
\]
Use the $6n$ version of HD-GC03. Its finite modes have exceptional
length at most $252$; mode one moves at most twelve letters;
mode infinity has index at least $6n/7-72$. Each large branch
therefore has index at least $n/2-252$. Five would exceed the
total natural index $2n+2g-2\leq2n+26$, so there are at most four.
The small total index satisfies
\[
                (n-12)J\leq A\leq4n+76,
                \qquad J<5,
\]
hence there are at most four small genuine branches. This proves
the eight-branch bound without classifying their exact types.

For an actual tail-two source satisfying (FT2P.1), HD-FT14 gives
$b\leq14$. Therefore at least four of its at least twelve genuine
original inertia operators project to the identity in $A_n$ or
$S_n$. The effective support estimate (FT2P.18), whose proof uses
only $b\geq0$, leaves total normalized moving rank at most $61/n^2$
for those entries. If $B$ denotes the tensor-partner dimension,
each such genuine entry has normalized moving rank at least $1/B$, so
\[
             B\geq4n^2/61,
 \qquad
 0<\frac{\operatorname{rank}(1-T_i)}{\dim V}
                       \leq\frac{61}{4n^2}
 \quad\hbox{for at least one partner-only entry}.  \tag{EG12}
\]
This is weaker than the seven-entry consequence above an unknown
$N_{14}$, but it is fully numerical. The precise pair-genus
classification cutoff, an unbounded-partner exclusion, and source
descent or detector construction are not asserted here.

\paragraph{Project frame.}
These effective inequalities constrain actual curve suppliers and
their weight-one cohomology. Their use in the earlier integral
projector and cycle-transfer results retains the exact condition-A,
condition-I and spreading locators already recorded there. They
do not turn a support inequality into an integral detecting cycle,
do not bound arbitrary original overgroups, and invoke none of
Markman, Severn or032. The determinant and Hodge-family gates
remain as stated until an actual source and its cycles are supplied.
